\title{Simplified State Space Layers for Sequence Modeling}

\begin{abstract}
Models using \emph{structured state space sequence} (S4) layers have achieved state-of-the-art performance on long-range sequence modeling tasks. An S4 layer combines linear state space models (SSMs), the HiPPO framework, and deep learning to achieve high performance. We build on the design of the S4 layer and introduce a new state space layer, the \emph{S5 layer}.  Whereas an S4 layer uses many independent single-input, single-output SSMs, the S5 layer uses one multi-input, multi-output SSM.  We establish a connection between S5 and S4, and use this to develop the initialization and parameterization used by the S5 model.  The result is a state space layer that can leverage efficient and widely implemented parallel scans, allowing S5 to match the computational efficiency of S4, while also achieving state-of-the-art performance on several long-range sequence modeling tasks.  S5 averages $87.4\%$ on the long range arena benchmark, and $98.5\%$ on the most difficult Path-X task.
\end{abstract}

\section{Introduction}
\label{sec:intro}

Efficiently modeling long sequences is a challenging problem in machine learning.  Information crucial to solving tasks may be encoded jointly between observations that are thousands of timesteps apart. Specialized variants of recurrent neural networks (RNNs)~\citep{arjovsky2016unitary, erichson2021lipschitz, rusch2021unicornn, chang2019antisymmetricrnn}, convolutional neural networks (CNNs)~\citep{bai2018empirical, oord2016wavenet, romero2021ckconv}, and transformers~\citep{vaswani2017attention} have been developed to try to address this problem. In particular, many efficient transformer methods have been introduced~\citep{choromanski2021rethinking, katharopoulos2020transformers, Kitaev2020Reformer:, beltagy2020longformer, gupta2020gmat, wang2020linformer} to address the standard transformer's quadratic complexity in the sequence length.  However, these more efficient transformers still perform poorly on very long-range sequence tasks~\citep{tay2020lra}. 

\citet{gu2021s4} presented an alternative approach using \emph{structured state space sequence} (S4) layers.  An S4 layer defines a nonlinear sequence-to-sequence transformation via a bank of many independent single-input, single-output (SISO) linear state space models (SSMs)~\citep{gu2021lssl}, coupled together with nonlinear mixing layers.  Each SSM leverages the HiPPO framework~\citep{gu2020hippo} by initializing with specially constructed state matrices.  Since the SSMs are linear, each layer can be equivalently implemented as a convolution, which can then be applied efficiently by parallelizing across the sequence length. Multiple S4 layers can be stacked to create a deep sequence model.  Such models have achieved significant improvements over previous methods, including on the \emph{long range arena} (LRA)~\citep{tay2020lra} benchmarks specifically designed to stress test long-range sequence models. Extensions have shown good performance on raw audio generation~\citep{goel2022sashimi} and classification of long movie clips~\citep{islam2022long}. 

We introduce a new state space layer that builds on the S4 layer, the \emph{S5} layer, illustrated in Figure~\ref{fig:banner:s5}.  S5 streamlines the S4 layer in two main ways.  First, S5 uses one multi-input, multi-output (MIMO) SSM in place of the bank of many independent SISO SSMs in S4.  Second, S5 uses an efficient and widely implemented parallel scan.  This removes the need for the convolutional and frequency-domain approach used by S4, which requires a non-trivial computation of the convolution kernel. The resulting state space layer has the same computational complexity as S4, but operates purely recurrently and in the time domain. 

We then establish a mathematical relationship between S4 and S5. This connection allows us to inherit the HiPPO initialization schemes that are key to the success of S4.  Unfortunately, the specific HiPPO matrix that S4 uses for initialization cannot be diagonalized in a numerically stable manner for use in S5. However, in line with recent work on the \emph{DSS}~\citep{gupta2022dss} and \emph{S4D}~\citep{gu2022parameterization} layers, we found that a diagonal approximation to the HiPPO matrix achieves comparable performance. We extend a result from~\citet{gu2022parameterization} to the MIMO setting, which justifies the diagonal approximation for use in S5. We leverage the mathematical relationship between S4 and S5 to inform several other aspects of parameterization and initialization, and we perform thorough ablation studies to explore these design choices.

The final S5 layer has many desirable properties. It is straightforward to implement (see Appendix~\ref{app:sec:implementation}),\footnote{The full S5 implementation is available at: \url{https://github.com/lindermanlab/S5}.} enjoys linear complexity in the sequence length, and can efficiently handle time-varying SSMs and irregularly sampled observations (which is intractable with the convolution implementation of S4).  S5 achieves state-of-the-art performance on a variety of long-range sequence modeling tasks, with an LRA average of ${87.4\%}$, and ${98.5\%}$ accuracy on the most difficult Path-X task. 

\section{Background}
\label{genbackground}
We provide the necessary background in this section prior to introducing the S5 layer in Section~\ref{sec:simplifying}. 

\subsection{Linear State Space Models}
\label{subsec:ssm}
Continuous-time linear SSMs are the core component of both the S4 layer and the S5 layer.  Given an input signal $\mathbf{u}(t)\in \mathbb{R}^U$, a latent state $\mathbf{x}(t)\in \mathbb{R}^P$ and an output signal $\mathbf{y}(t)\in \mathbb{R}^{M}$, a linear continuous-time SSM is defined by the differential equation:

\begin{align}
    \dfrac{\mathrm{d} \mathbf{x}(t)}{\mathrm{d} t} = \mathbf{A}\mathbf{x}(t) + \mathbf{B}\mathbf{u}(t), \quad \quad \mathbf{y}(t) = \mathbf{C}\mathbf{x}(t) + \mathbf{D}\mathbf{u}(t) , \label{eq:cont_ssm} 
\end{align}

and is parameterized by a state matrix $\mathbf{A}\in \mathbb{R}^{P \times P}$, an input matrix $\mathbf{B}\in \mathbb{R}^{P \times U}$, an output matrix $\mathbf{C}\in \mathbb{R}^{M \times P}$ and a feedthrough matrix $\mathbf{D}\in \mathbb{R}^{M \times U}$.  For a constant step size, $\Delta$, the SSM can be \emph{discretized} using, e.g. Euler, bilinear or zero-order hold (ZOH) methods to define the linear recurrence

\begin{align}
    \mathbf{x}_k = \overline{\mathbf{A}}\mathbf{x}_{k-1} + \overline{\mathbf{B}}\mathbf{u}_k, \quad \quad \mathbf{y}_k = \overline{\mathbf{C}}\mathbf{x}_k + \overline{\mathbf{D}}\mathbf{u}_k, \label{eq:discrete_ssm}
\end{align}

where the discrete-time parameters are each a function, specified by the discretization method, of the continuous-time parameters.  See \citet{iserles2009first} for more information on discretization methods.  

\subsection{Parallelizing Linear State Space Models With Scans}
\label{subsec:parscan}
We use parallel scans to efficiently compute the states of a discretized linear SSM.  Given a binary associative operator $\bullet$ (i.e. $(a \bullet b) \bullet c = a \bullet (b \bullet c)$) and a sequence of $L$ elements $[a_1, a_2, ..., a_{L}]$, the scan operation (sometimes referred to as \emph{all-prefix-sum}) returns the sequence

\begin{align}
    [a_1, \ (a_1 \bullet a_2), \ ..., \ (a_1 \bullet a_2 \bullet ... \bullet a_{L})].
\end{align}

Computing a length $L$ linear recurrence of a discretized SSM, $\mathbf{x}_k = \overline{\mathbf{A}}\mathbf{x}_{k-1} + \overline{\mathbf{B}}\mathbf{u}_k$ as in (\ref{eq:discrete_ssm}), is a specific example of a scan operation.  As discussed in  Section 1.4 of \citet{blelloch1990prefix}, parallelizing the linear recurrence of the latent transitions in the discretized SSM above can be computed in a parallel time of $\mathcal{O}(T_{\odot}\log L)$, assuming $L$ processors, where $T_{\odot}$ represents the cost of matrix-matrix multiplication. For a general matrix $\overline{\mathbf{A}}\in \mathbb{R}^{P \times P}$, $T_\odot$ is $\mathcal{O}(P^3)$.  This can be prohibitively expensive in deep learning settings. However, if $\overline{\mathbf{A}}$ is a diagonal matrix, the parallel time becomes $\mathcal{O}(P \log L)$ with $L$ processors and only requires $\mathcal{O}(PL)$ space. Finally, we note that efficient parallel scans are implemented in a work-efficient manner, thus the total computational cost of the parallel scan with a diagonal matrix is $\mathcal{O}(PL)$ operations. See Appendix \ref{app:sec:parallel_scan} for more information on parallel scans.

\subsection{S4: Structured State Space Sequence Layers}
\label{background:S4}
The S4 layer \citep{gu2021s4} defines a nonlinear sequence-to-sequence transformation, mapping from an input sequence $\mathbf{u}_{1:L} \in \mathbb{R}^{L \times H}$ to an output sequence $\mathbf{u}'_{1:L} \in \mathbb{R}^{L \times H}$.  An S4 layer contains a bank of $H$ independent single-input, single-output (SISO) SSMs with $N$-dimensional states. Each \emph{S4 SSM} is applied to one dimension of the input sequence.  This results in an independent linear transformation from each input channel to each preactivation channel.  A nonlinear activation function is then applied to the preactivations.  Finally, a position-wise linear \emph{mixing} layer is applied to combine the independent features and produce the output sequence $\mathbf{u}'_{1:L}$.  Figure \ref{app:fig:banner:s4} in the appendix illustrates the view of the S4 layer as a bank of independent SSMs. Figure \ref{fig:matrix:s4} shows an alternative view of S4 as one large SSM with state size $HN$ and block-diagonal state, input and output matrices.  

Each S4 SSM leverages the HiPPO framework for online function approximation~\citep{gu2020hippo} by initializing the state matrices with a HiPPO matrix (most often the HiPPO-LegS matrix).  This was demonstrated empirically to lead to strong performance~\citep{gu2021lssl, gu2021s4}, and can be shown as approximating long-range dependencies with respect to an infinitely long, exponentially-decaying measure~\citep{gu2022train}.  While the HiPPO-LegS matrix is not stably diagonalizable~\citep{gu2021s4}, it can be represented as a normal plus low-rank (NPLR) matrix.   The normal component, referred to as HiPPO-N and denoted $\mathbf{A}_{\mathrm{LegS}}^{\mathrm{Normal}}$, can be diagonalized.  Thus, the HiPPO-LegS can be conjugated into a diagonal plus low-rank (DPLR) form, which S4 then utilizes to derive an efficient form of the convolution kernel.  This motivates S4's DPLR parameterization.

Efficiently applying the S4 layer requires two separate implementations depending on context: a recurrent mode and a convolution mode.  For online generation, the SSM is iterated recurrently, much like other RNNs.  However, when the entire sequence is available and the observations are evenly spaced, a more efficient convolution mode is used.  This takes advantage of the ability to represent the linear recurrence as a one-dimensional convolution between the inputs and a convolution kernel for each of the SSMs.  Fast Fourier transforms (FFTs) can then be applied to efficiently parallelize this application.  Figure \ref{app:fig:banner:s4} in the appendix illustrates the convolution approach of the S4 layer for offline processing.  We note that while parallel scans could, in principle, allow a recurrent approach to be used in offline scenarios, applying the parallel scan to all $H$ of the $N$-dimensional SSMs would in general be much more expensive than the convolution approach S4 actually uses.

The trainable parameters of each S4 layer are the $H$ independent copies of the learnable SSM parameters and the $\mathcal{O}(H^2)$ parameters of the mixing layer. For each of the $h\in \{1,...,H\}$ S4 SSMs, given a scalar input signal $u^{(h)}(t)\in \mathbb{R}$, an S4 SSM uses an input matrix $\mathbf{B}^{(h)}\in \mathbb{C}^{N\times 1}$, a DPLR parameterized transition matrix $\mathbf{A}^{(h)}\in \mathbb{C}^{N \times N}$, an output matrix $\mathbf{C}^{(h)}\in \mathbb{C}^{1\times N}$, and feedthrough matrix $\mathbf{D}^{(h)}\in \mathbb{R}^{1\times 1}$, to produce a signal $y^{(h)}(t) \in \mathbb{R}$. To apply the S4 SSMs to discrete sequences, each continuous-time SSM is discretized using a constant timescale parameter $\Delta^{(h)}\in \mathbb{R}_+$. The learnable parameters of each SSM are the timescale parameter $\Delta^{(h)}\in \mathbb{R}_+$, the continuous-time parameters $\mathbf{B}^{(h)}$, $\mathbf{C}^{(h)}$, $\mathbf{D}^{(h)}$, and the DPLR matrix, parameterized by vectors $\mathbf{\Lambda}^{(h)}\in\mathbb{C}^{N}$ and $\mathbf{p}^{(h)}, \mathbf{q}^{(h)} \in \mathbb{C}^{N}$ representing the diagonal matrix and low-rank terms respectively.  For notational compactness we denote the concatenation of the $H$ S4 SSM states at discrete time index $k$ as $\mathbf{x}_k^{(1:H)} = \big[ (\mathbf{x}_k^{(1)})^\top , \ldots , (\mathbf{x}_k^{(H)})^\top \big]^\top$, and the $H$ SSM outputs as $\mathbf{y}_k = \big[ \mathbf{y}_k^{(1)} , \ldots , \mathbf{y}_k^{(H)} \big]^\top$.

\renewcommand{.8}{.8}

\section{The S5 Layer}
\label{sec:simplifying}
In this section we present the S5 layer.  We describe its structure, parameterization and computation, particularly focusing on how each of these differ from S4. 

\subsection{S5 Structure: From SISO to MIMO}
\label{subsec:simplestruct}
The S5 layer replaces the bank of SISO SSMs (or large block-diagonal system) in S4 with a multi-input, multi-output (MIMO) SSM, as in (\ref{eq:cont_ssm}), with a latent state size $P$, and input and output dimension $H$. The discretized version of this MIMO SSM can be applied to a vector-valued input sequence $\mathbf{u}_{1:L}\in \mathbb{R}^{L\times H}$, to produce a vector-valued sequence of \emph{SSM outputs} (or preactivations) $\mathbf{y}_{1:L}\in \mathbb{R}^{L\times H}$, using latent states $\mathbf{x}_k\in \mathbb{R}^{P}$.  A nonlinear activation function is then applied to produce a sequence of layer outputs $\mathbf{u}'_{1:L}\in \mathbb{R}^{L\times H}$. See Figure \ref{fig:matrix:S5} for an illustration. Unlike S4, we do not require an additional position-wise linear layer, since these features are already mixed.  We note here that compared to the $HN$ latent size of the block-diagonal SSM in the S4 layer, S5's latent size $P$ can be significantly smaller, allowing for the use of efficient parallel scans, as we discuss in Section \ref{subsec:simplecomp}. 

\subsection{S5 Parameterization: Diagonalized Dynamics}
\label{subsec:simpleparam}

The parameterization of the S5 layer's MIMO SSM is motivated by the desire to use efficient parallel scans. As discussed in Section \ref{subsec:parscan}, a diagonal state matrix is required to efficiently compute the linear recurrence using a parallel scan. Thus, we diagonalize the system, writing the continuous-time state matrix as  $\mathbf{A}=\mathbf{V}\mathbf{\Lambda}\mathbf{V}^{-1}$, where $\mathbf{\Lambda}\in \mathbb{C}^{P \times P}$ denotes the diagonal matrix containing the eigenvalues and $\mathbf{V} \in \mathbb{C}^{P \times P}$ corresponds to the eigenvectors.  Therefore, we can diagonalize the continuous-time latent dynamics from (\ref{eq:cont_ssm}) as

\begin{align}
    \dfrac{\mathrm{d} \mathbf{V}^{-1}\mathbf{x}(t)}{\mathrm{d} t} = \boldsymbol\Lambda \mathbf{V}^{-1}\mathbf{x}(t) + \mathbf{V}^{-1}\mathbf{B}\mathbf{u}(t).
    \label{eq:S5eigenbasis}
\end{align}

Defining $\tilde{\mathbf{x}}(t)= \mathbf{V}^{-1}\mathbf{x}(t)$, $\tilde{\mathbf{B}}= \mathbf{V}^{-1}\mathbf{B}$, and  $\Tilde{\mathbf{C}}=\mathbf{C}\mathbf{V}$ gives a reparameterized system,

\begin{align}
    \dfrac{\mathrm{d} \tilde{\mathbf{x}}(t)}{\mathrm{d} t} = \mathbf{\Lambda} \tilde{\mathbf{x}}(t) + \tilde{\mathbf{B}}\mathbf{u}(t), \quad \quad \mathbf{y}(t) &= \tilde{\mathbf{C}}\tilde{\mathbf{x}}(t) +  \mathbf{D}\mathbf{u}(t).
    \label{eq:S5ssm} 
\end{align}

This is a linear SSM with a diagonal state matrix. This diagonalized system can be discretized with a timescale parameter $\Delta\in \mathbb{R}_+$ using the ZOH method to give another diagonalized system with parameters

\begin{align}
    \label{eq:diag_disc_params}
    \overline{\mathbf{\Lambda}} = e^{\mathbf{\Lambda}\Delta}, \quad \overline{\mathbf{B}} = \mathbf{\Lambda}^{-1}(\overline{\mathbf{\Lambda}}-\mathbf{I})\tilde{\mathbf{B}}, \quad \overline{\mathbf{C}} = \tilde{\mathbf{C}}, \quad \overline{\mathbf{D}} = \mathbf{D}. 
\end{align}

In practice, we use a vector of learnable timescale parameters $\mathbf{\Delta}\in \mathbb{R}^P$ (see Section \ref{sec:s4s5rel:untying}) and restrict the feedthrough matrix $\mathbf{D}$ to be diagonal.  The S5 layer therefore has the learnable parameters:  $\tilde{\mathbf{B}}\in \mathbb{C}^{P \times H}$, $\tilde{\mathbf{C}}\in \mathbb{C}^{H \times P}$, $\mathrm{diag}(\mathbf{D})\in \mathbb{R}^{H}$,  $\mathrm{diag}(\mathbf{\Lambda}) \in \mathbb{C}^{P}$, and $\mathbf{\Delta}\in \mathbb{R}^P$. 

\paragraph{Initialization}
Prior work showed that the performance of deep state space models are sensitive to the initialization of the state matrix~\citep{gu2021lssl, gu2021s4}.  We discussed in Section \ref{subsec:parscan} that state matrices must be diagonal for efficient application of parallel scans.  We also discussed in Section \ref{background:S4} that the HiPPO-LegS matrix cannot be diagonalized stably, but that the HiPPO-N matrix can be.  In Section \ref{sec:s4s5rel} we connect the dynamics of S5 to S4 to suggest why initializing with HiPPO-like matrices may also work well in the MIMO setting.  We support this empirically, finding that diagonalizing the HiPPO-N matrix leads to good performance, and perform ablations in Appendix \ref{app:section:ablations} to compare to other initializations.  We note that DSS~\citep{gupta2022dss} and S4D~\citep{gu2022parameterization} layers also found strong performance in the SISO setting by using a diagonalization of the HiPPO-N matrix.

\paragraph{Conjugate Symmetry} The complex eigenvalues of a diagonalizable matrix with real entries always occur in conjugate pairs.  We enforce this conjugate symmetry by using half the number of eigenvalues and latent states. This ensures real outputs and reduces the runtime and memory usage of the parallel scan by a factor of two.  This idea is also discussed in \citet{gu2022parameterization}. 

\subsection{S5 Computation: Fully Recurrent} 
\label{subsec:simplecomp}
Compared to the large $HN$ effective latent size of the block-diagonal S4 layer, the smaller latent dimension of the S5 layer ($P$) allows the use of efficient parallel scans when the entire sequence is available.  The S5 layer can therefore be efficiently used as a recurrence in the time domain for both online generation and offline processing.  Parallel scans and the continuous-time parameterization also allow for efficient handling of irregularly sampled time series and other time-varying SSMs, by simply supplying a different $\overline{\mathbf{A}}_k$ matrix at each step.  We leverage this feature and apply S5 to irregularly sampled data in Section \ref{sec:exp:variable}.  In contrast, the convolution of the S4 layer requires a time invariant system and regularly spaced observations.

\subsection{Matching the Computational Efficiency of S4 and S5} A key design desiderata for S5 was matching the computational complexity of S4 for both online generation and offline recurrence. The following proposition guarantees that their complexities are of the same order if S5's latent size $P=\mathcal{O}(H)$.

\begin{prop} 
\label{prop:efficiency}
Given an S4 layer with $H$ input/output features, an S5 layer with $H$ input/output features and a latent size $P=\mathcal{O}(H)$ has the same order of magnitude complexity as an S4 layer in terms of both runtime and memory usage. 
\end{prop}

\begin{proof}
See Appendix \ref{app:comp_proof}.
\end{proof}

We also support this proposition with empirical comparisons in Appendix \ref{apps:runtimecomp}.

\section{Relationship Between S4 and S5}
\label{sec:s4s5rel}
We now establish a relationship between the dynamics of S5 and S4.  In Section \ref{app:sec:equiv} we show that, under certain conditions, the outputs of the S5 SSM can be \emph{interpreted} as a projection of the latent states computed by a particular S4 system. This interpretation motivates using HiPPO initializations for S5, which we discuss in more detail in Section \ref{sec:s4s5rel:diag}.  In Section \ref{sec:s4s5rel:untying} we discuss how the conditions required to relate the dynamics further motivate initialization and parameterization choices.  

\subsection{Different Output Projections of Equivalent Dynamics}
\label{app:sec:equiv}
We compare the dynamics of S4 and S5 under some simplifying assumptions:

\begin{assumption}
\label{ass:h_to_h}
We consider only $H$-dimensional to $H$-dimensional sequence maps.
\end{assumption}

\begin{assumption}
\label{ass:tied_a}
We assume the state matrix of each S4 SSM is identical, $\mathbf{A}^{(h)} = \mathbf{A} \in \mathbb{C}^{N \times N}$.
\end{assumption}

\begin{assumption}
\label{ass:tied_delta}
We assume the timescales of each S4 SSM are identical, $\Delta^{(h)} = \Delta \in \mathbb{R}_+$
\end{assumption}

\begin{assumption}
\label{ass:tied_s5}
We assume that the same state matrix $\mathbf{A}$ is used in S5 as in S4 (also cf. Assumption \ref{ass:tied_a}).  Note this also specifies the S5 latent size $P=N$.  We also assume the S5 input matrix is the horizontal concatenation of the column input vectors used by S4: $\mathbf{B} \triangleq \left[ \mathbf{B}^{(1)} \mid \ldots \mid \mathbf{B}^{(H)} \right]$.
\end{assumption}

We will discuss relaxing these assumptions shortly, but under these conditions it is straightforward to derive a relationship between the dynamics of S4 and S5:

\begin{prop}
\label{prop:equiv}
Consider an S5 layer, with state matrix $\mathbf{A}$, input matrix $\mathbf{B}$ and some output matrix $\mathbf{C}$ (cf. Assumption \ref{app:ass:h_to_h}); and an S4 layer, where each of the $H$ S4 SSMs has state matrix $\mathbf{A}$ (cf. Assumption \ref{app:ass:tied_a}, \ref{ass:tied_s5}) and input vector $\mathbf{B}^{(h)}$ (cf. Assumption \ref{app:ass:tied_s5}).  If the S4 and S5 layers are discretized with the same timescales (cf. Assumption \ref{ass:tied_delta}), then the S5 SSM produces outputs, $\mathbf{y}_k$, equivalent to a linear combination of the latent states of the $H$ S4 SSMs, $\mathbf{y}_k = \mathbf{C}^{\mathrm{equiv}}\mathbf{x}_k^{(1:H)}$, where\ $\mathbf{C}^{\mathrm{equiv}} = \left[\ \mathbf{C}\ \cdots\ \mathbf{C}\ \right]$.
\end{prop}

\begin{proof}
See Appendix \ref{app:sec:relation:equiv}.
\end{proof}

Importantly, the S5 SSM outputs are \emph{not} equal to the outputs of the block-diagonal S4 SSM.  Instead they are equivalent to the outputs of the block-diagonal S4 SSM with modified output matrix $\mathbf{C}^{\mathrm{equiv}}$.  Under the assumptions, however, the underlying state dynamics \emph{are} equivalent.  Recalling that initializing the S4 dynamics with HiPPO was key to performance~\citep{gu2021s4}, the relationship established in Proposition~\ref{prop:equiv} motivates using HiPPO initializations for S5, as we now discuss.  

\subsection{Diagonalizable Initialization}
\label{sec:s4s5rel:diag}
Ideally, given the interpretation above, we would initialize S5 with the exact HiPPO-LegS matrix. Unfortunately, as discussed in Section \ref{background:S4}, this matrix is not stably diagonalizable, as is required for the efficient parallel scans used for S5. However, \citet{gupta2022dss} and \citet{gu2022parameterization} showed empirically that removing the low rank terms and initializing with the diagonalized HiPPO-N matrix still performed well. \citet{gu2022parameterization} offered a theoretical justification for the use of this normal approximation for single-input systems: in the limit of infinite state dimension, the linear ODE with HiPPO-N state matrix produces the same dynamics as an ODE with the HiPPO-LegS matrix.  Using linearity, it is straightforward to extend this result to the multi-input system that S5 uses:

\begin{corollary}[Extension of Theorem 3 in \citet{gu2022parameterization}]
\label{prop:s4d_extension}
Consider $\mathbf{A}_{\mathrm{LegS}}\in \mathbb{R}^{N\times N}$, $\mathbf{A}_{\mathrm{LegS}}^{\mathrm{Normal}}\in \mathbb{R}^{N\times N}$, $\mathbf{B}_{\mathrm{LegS}}\in \mathbb{R}^{N \times H}, \mathbf{P}_{\mathrm{LegS}}\in \mathbb{R}^{N}$ as defined in Appendix \ref{subsec:initA}. Given vector-valued inputs $\mathbf{u}(t)\in \mathbb{R}^H$, the ordinary differential equation $\dfrac{\mathrm{d} \mathbf{x}'(t)}{\mathrm{d} t} = \mathbf{A}_{\mathrm{LegS}}^{\mathrm{Normal}} \mathbf{x}'(t) + \frac{1}{2}\mathbf{B}_{\mathrm{LegS}}\mathbf{u}(t)$ converges to $\dfrac{\mathrm{d} \mathbf{x}(t)}{\mathrm{d} t} = \mathbf{A}_{\mathrm{LegS}}\mathbf{x}(t) + \mathbf{B}_{\mathrm{LegS}}\mathbf{u}(t)$ as $N \rightarrow \infty$.
\end{corollary}

We include a simple proof of this extension in Appendix \ref{app:sec:relation:init}.  This extension motivates the use of HiPPO-N to initialize S5's MIMO SSM. Note that S4D (the diagonal extension of S4) uses the same HiPPO-N matrix. Thus, when under the assumptions in Proposition \ref{prop:equiv}, an S5 SSM in fact produces outputs that are equivalent to a linear combination of the latent states produced by S4D's SSMs. Our empirical results in Section \ref{sec:exp} suggest that S5 initialized with the HiPPO-N matrix performs just as well as S4 initialized with the HiPPO-LegS matrix.  

\subsection{Relaxing the Assumptions}
\label{sec:s4s5rel:untying}

We now revisit the assumptions required for Proposition \ref{prop:equiv}, since they only relate a constrained version of S5 to a constrained version of S4. Regarding Assumption \ref{ass:tied_a}, \citet{gu2021s4} report that S4 models with tied state matrices can still perform well, though allowing different state matrices often yields higher performance. Likewise, requiring a single scalar timescale across all of the S4 SSMs, per Assumption \ref{ass:tied_delta}, is restrictive.  S4 typically learns different timescale parameters for each SSM~\citep{gu2022train} to capture different timescales in the data. To relax these assumptions, note that Assumption \ref{ass:tied_s5} constrains S5 to have dimension $P=N$, and $N$ is typically much smaller than the dimensionality of the inputs, $H$. Proposition~\ref{prop:efficiency} established that S5 can match S4's complexity with $P = \mathcal{O}(H)$. By allowing for larger latent state sizes, Assumptions 2 and 3 can be relaxed, as discussed in Appendix \ref{app:sec:relation:blocks}. We also discuss how this relaxation motivates a block-diagonal initialization with HiPPO-N matrices on the diagonal. Finally, to further relax the tied timescale assumptions, we note that in practice, we find improved performance by learning $P$ different timescales (one per state). See Appendix \ref{app:sec:relation:timescales} for further discussion of this empirical finding and the ablations in Appendix \ref{app:simo-ablation}.

\section{Related work}
\label{sec: related_worrk} S5 is most directly related to S4 and its other extensions, which we have discussed thoroughly. However, there is prior literature that uses similar ideas to those developed here. For example, prior work studied approximating nonlinear RNNs with stacks of linear RNNs connected by nonlinear layers, while also using parallel scans~\citep{martin2018parallelizing}.  \citet{martin2018parallelizing} showed that several efficient RNNs, such as QRNNs~\citep{bradbury2016quasi} and SRUs~\citep{lei2017simple}, fall into a class of linear surrogate RNNs that can leverage parallel scans. \citet{kaul2020linear} also used parallel scans for an approach that approximates RNNs with stacks of discrete-time single-input, multi-output (SIMO) SSMs.  However, S4 and S5 are the only methods to significantly outperform other comparable state-of-the-art nonlinear RNNs, transformers and convolution approaches.  Our ablation study in Appendix \ref{app:ablation:disc_hippo} suggests that this performance gain over prior attempts at parallelized linear RNNs is likely due to the continuous-time parameterization and the HiPPO initialization.

\section{Experiments}
\label{sec:exp}
We now compare empirically the performance of the S5 layer to the S4 layer and other baseline methods.  We use the S5 layer as a drop-in replacement for the S4 layer.  The architecture consists of a linear input encoder, stacks of S5 layers, and a linear output decoder~\citep{gu2021s4}.  For all experiments we choose the S5 dimensions to ensure similar computational complexities as S4, following the conditions discussed in Section \ref{subsec:simplecomp}, as well as comparable parameter counts.  The results we present show that the S5 layer matches the performance and efficiency of the S4 layer.  We include in the appendix further ablations, baselines and runtime comparisons.

\subsection{Long Range Arena}
\label{subsec:LRA}
The long range arena (LRA) benchmark~\citep{tay2020lra} is a suite of six sequence modeling tasks, with sequence lengths from 1,024 to over 16,000.  The suite was specifically developed to benchmark the performance of architectures on long-range modeling tasks (see Appendix \ref{app:sec:configs} for more details). Table \ref{tab:lra} presents S5's LRA performance in comparison to other methods. S5 achieves the highest average score among methods that have linear complexity in sequence length (most notably S4, S4D, and the concurrent works: Liquid-S4~\citep{hasani2022liquid} and Mega-chunk~\citep{ma2022mega}). Most significantly, S5 achieves the highest score among all models (including Mega~\citep{ma2022mega}) on the Path-X task, which has by far the longest sequence length of the tasks in the benchmark.  

\subsection{Raw Speech Classification}
\label{subsec:speech}
The Speech Commands dataset~\citep{warden2018speech} contains high-fidelity sound recordings of different human readers reciting a word from a vocabulary of 35 words.  The task is to classify which word was spoken.  We show in Table \ref{tab:sc35_compact} that S5 outperforms the baselines, outperforms previous S4 methods and performs similarly to the concurrent Liquid-S4 method~\citep{hasani2022liquid}. As S4 and S5 methods are parameterized in continuous-time, these models can be applied to datasets with different sampling rates without the need for re-training, simply by globally re-scaling the timescale parameter $\Delta$ by the ratio between the new and old sampling rates.  The result of applying the best S5 model \emph{trained on 16kHz data}, to the speech data sampled (via decimation) at 8kHz, without any additional fine-tuning, is also presented in Table \ref{tab:sc35_compact}.  S5 also improves this metric over the baseline methods. 

\subsection{Variable Observation Interval}
\label{sec:exp:variable}
The final application we study here highlights how S5 can naturally handle observations received at irregular intervals.  S5 does so by supplying a different $\Delta_t$ value to the discretization at each step.  We use the pendulum regression example presented by \citet{becker2019recurrent} and \citet{schirmer2022modeling}, illustrated in Figure \ref{fig:exp:pendulum_illustration}.  The input sequence is a sequence of $L=50$ images, each $24 \times 24$ pixels in size, that has been corrupted with a correlated noise process and sampled at irregular intervals from a continuous trajectory of duration $T=100$.  The targets are the sine and cosine of the angle of the pendulum, which follows a nonlinear dynamical system.  The velocity is unobserved.  We match the architecture, parameter count and training procedure of \citet{schirmer2022modeling}.  Table \ref{tab:results:pendulum_mse} summarizes the results of this experiment.  S5 outperforms CRU on the regression task, recovering a lower mean error.  Furthermore, S5 is markedly faster than CRU on the same hardware.  

\subsection{Pixel-level 1-D Image Classification}
\label{subsec:pixel}
Table~\ref{tab:pixellong} in Appendix \ref{sec:pixellong} shows results of S5 on other common benchmarks including sequential MNIST, permuted sequential MNIST and sequential CIFAR (color).  We see that S5 broadly matches the performance of S4, and outperforms a range of state-of-the-art RNN-based methods.  

\section{Conclusion}
\label{sec:conclusion}
We introduce the S5 layer for long-range sequence modeling.  The S5 layer modifies the internal structure of the S4 layer, and replaces the frequency-domain approach used by S4 with a purely recurrent, time-domain approach leveraging parallel scans. S5 achieves high performance while retaining the computational efficiency of S4.  S5 also provides further opportunities.  For instance, unlike the convolutional S4 methods, the parallel scan unlocks the ability to efficiently and easily process time-varying SSMs whose parameters can vary with time. Section \ref{sec:exp:variable} illustrated an example of this for sequences sampled at variable sampling rates. The concurrently developed method, \emph{Liquid-S4}~\citep{hasani2022liquid}, uses an input-dependent bilinear dynamical system and highlights further opportunities for time-varying SSMs. The more general MIMO SSM design will also enable connections to be made with classical probabilistic state space modeling as well as more recent work on parallelizing filtering and smoothing operations~\citep{sarkka2020temporal}.  More broadly, we hope the simplicity and generality of the S5 layer can expand the use of state space layers in deep sequence modeling and lead to new formulations and extensions.  

\end{document}
